\documentclass[10pt,twocolumn]{article}

\usepackage[margin=1in]{geometry}
\usepackage{amsmath,amssymb}
\usepackage{graphicx}
\usepackage{booktabs}
\usepackage{multirow}
\usepackage{xcolor}
\usepackage{hyperref}
\usepackage{natbib}
\usepackage{tablefootnote}

\definecolor{cliffred}{RGB}{204,51,51}
\definecolor{safegreen}{RGB}{51,153,51}

\title{How Vision Connectors Shape Token Pruning Robustness\\ in Vision-Language-Action Models}

\author{Anonymous}
\date{}

\begin{document}
\maketitle

%==============================================================================
% ABSTRACT
%==============================================================================
\begin{abstract}
Inference-time vision token pruning is a promising approach to accelerate Vision-Language-Action (VLA) models, yet its effectiveness varies dramatically across architectures for reasons that remain poorly understood.
We present the first systematic cross-architecture analysis of token pruning robustness in VLAs, evaluating 5 models (450M--7B parameters) across 3 connector families and 4 task suites.
We find that pruning robustness is primarily explained by \emph{connector architecture}---whether the vision connector produces token-independent or token-entangling representations---with task complexity acting as a modulating factor.
Token importance Coefficient of Variation (CV), a zero-cost single-forward-pass diagnostic, predicts both pruning tolerance and optimal pruning strategy: low-CV models ($\text{CV}\!\leq\!0.15$) are actively harmed by importance-based pruning at moderate compression ($k\!=\!0.5$), with random pruning outperforming L2-norm selection by up to 36 percentage points (pp).
Pixel-shuffle connectors exhibit catastrophic failure---removing just 5\% of tokens yields 0\% task success---while linear projections tolerate up to 90\% removal.
For the worst-case pixel-shuffle architecture, we propose Vision Token Re-encoding (VTR), a lightweight 22M-parameter module that recovers from 0\% (hard pruning) to near-baseline performance (60.1\% vs.\ 59.2\% baseline) at any compression ratio from 0\% to 87.5\%.
VTR is ineffective on MLP-projector architectures (0\% success), confirming that optimal compression strategy is fundamentally connector-dependent.
\end{abstract}

%==============================================================================
% 1. INTRODUCTION
%==============================================================================
\section{Introduction}

Vision-Language-Action (VLA) models integrate pretrained vision encoders with large language models to generate robotic actions from visual observations and language instructions~\cite{brohan2023rt2,kim2024openvla,black2024pi0}.
A key deployment bottleneck is the quadratic cost of self-attention over vision tokens: models like $\pi$0.5 process 768 tokens per observation, consuming significant compute and memory in the KV-cache.

A growing body of work proposes inference-time token pruning to address this bottleneck~\cite{chen2024fastv,tang2024vlaiap,zhang2024lightvla}.
These methods report strong results on specific architectures, but an implicit assumption underlies them all: that token pruning is a \emph{general} technique whose effectiveness is primarily determined by the pruning algorithm.
We show this assumption is false.

Across 5 VLA models spanning 3 connector families (linear projection, multi-layer MLP, pixel-shuffle) and 4 model sizes (450M--7B), we find that \textbf{the vision connector is the primary determinant of pruning robustness}, not the pruning method.
Specifically, we identify a two-axis framework:

\noindent\textbf{Axis 1---Token Coupling.} Connectors that preserve token independence (linear projection, MLP) produce graded degradation under pruning.
Connectors that entangle tokens across spatial positions (BART encoder, pixel-shuffle) create cascading failure modes.
The importance CV of vision tokens, computable from a single forward pass, diagnoses a model's position on this axis ($\text{CV}=0.038$ for the most independent to $\text{CV}=0.314$ for the most entangled).

\noindent\textbf{Axis 2---Task Complexity.} For token-entangling connectors, task complexity modulates pruning sensitivity: X-VLA drops only 1.6\,pp on simple manipulation (Object suite) at $k\!=\!0.9$ but 28.2\,pp on multi-step reasoning (Goal suite) at the same keep ratio.
For token-independent connectors, this amplification is substantially weaker ($<$15\,pp for $\pi$0.5).

Our contributions are:
\begin{enumerate}
\item The first systematic cross-architecture, cross-task token pruning analysis in VLAs (5 models $\times$ 4 suites), revealing that connector architecture---not pruning algorithm---is the primary determinant of robustness.
\item Token importance CV as a zero-cost diagnostic that predicts both pruning tolerance and optimal pruning strategy, validated with Spearman $\rho=0.82$ ($p=0.03$) at moderate compression.
\item Vision Token Re-encoding (VTR), a targeted 22M-parameter recovery method for the worst-case pixel-shuffle architecture, achieving near-lossless compression at any ratio where hard deletion yields 0\%.
\end{enumerate}

%==============================================================================
% 2. RELATED WORK
%==============================================================================
\section{Related Work}

\paragraph{Token Pruning in Vision-Language Models.}
FastV~\cite{chen2024fastv} prunes vision tokens based on attention scores in VLMs, achieving 45\% FLOPs reduction with minimal quality loss.
VLA-IAP~\cite{tang2024vlaiap} extends this to VLAs, reporting 97.8\% success rate on OpenVLA under aggressive pruning.
LightVLA~\cite{zhang2024lightvla} and See What Matters~\cite{seewhatmatters2026} propose architecture-specific pruning strategies.
All of these methods are evaluated on MLP-projector VLAs.
Our work reveals why: MLP projectors produce independently meaningful tokens that tolerate pruning, while pixel-shuffle connectors create fundamentally different compressibility properties.

\paragraph{Vision-Language-Action Models.}
We study 5 VLA models spanning diverse design choices: $\pi$0~\cite{black2024pi0} and $\pi$0.5 use linear projections from PaliGemma; OpenVLA~\cite{kim2024openvla} and OpenVLA-OFT employ 3-layer MLP projectors; X-VLA processes tokens through a 12-layer BART encoder; SmolVLA uses pixel-shuffle compression.
This architectural diversity has not been previously studied in the context of token pruning.

%==============================================================================
% 3. EXPERIMENTAL SETUP
%==============================================================================
\section{Experimental Setup}

\subsection{Models and Connectors}

Table~\ref{tab:models} summarizes the 5 VLA models studied, spanning 3 connector families and an order of magnitude in parameter count.

\begin{table}[t]
\centering
\caption{VLA models studied. CV = Coefficient of Variation of token importance (L2-norm). Entropy = mean attention entropy across layers. Connector ``type'' refers to the effective vision-to-LLM pathway.}
\label{tab:models}
\small
\begin{tabular}{@{}lllrcc@{}}
\toprule
Model & Params & Connector & Tokens & CV & Entropy \\
\midrule
$\pi$0 & 3.3B & Linear proj. & 256 & 0.038 & 0.959 \\
$\pi$0.5 & 4B & Linear proj. & 768 & 0.055 & 0.666 \\
X-VLA & 880M & Linear + BART 12L & 50 & 0.151 & 0.614 \\
OpenVLA & 7B & 3-layer MLP & 256 & 0.215 & 0.661 \\
OFT & 7B & 3-layer MLP & 512 & 0.219 & 0.667 \\
SmolVLA & 450M & Pixel-shuffle & 128 & 0.314 & 0.952 \\
\bottomrule
\end{tabular}
\end{table}

\subsection{Benchmark and Evaluation}

We evaluate on \textbf{LIBERO}~\cite{liu2023libero}, using all 4 task suites: Object (single-step manipulation with visually similar objects), Spatial (spatial reasoning), Goal (multi-step goal-conditioned), and Long (long-horizon multi-step).
Each suite contains 10 tasks evaluated over 50 episodes per task (500 total).
Success rate (SR) is reported as the primary metric.

\subsection{Pruning Protocol}

We apply post-hoc token pruning at inference time without retraining.
For each model, we evaluate at keep ratios $k \in \{0.95, 0.9, 0.85, 0.8, 0.7, 0.5, 0.3\}$ using three pruning strategies:
(1)~\textbf{L2-norm}: retain tokens with highest L2-norm of the vision embedding;
(2)~\textbf{Random}: uniformly sample tokens without replacement;
(3)~\textbf{Attention-based}: retain tokens with highest aggregated attention weight.
Multi-seed evaluation (3 seeds) is used for key conditions.

\subsection{Diagnostic Metrics}

\textbf{Importance CV.} For each model, we compute the L2-norm of each vision token across a calibration set, then compute the coefficient of variation (standard deviation divided by mean) across tokens.
This requires a single forward pass and no labels.

\textbf{Attention Entropy.} Mean entropy of the attention distribution from the LLM's suffix tokens to vision tokens, averaged across layers.

%==============================================================================
% 4. RESULTS
%==============================================================================
\section{Results}

\subsection{Two-Axis Framework Overview}
\label{sec:framework}

Our central finding is that token pruning robustness is primarily explained by connector architecture, with task complexity as a modulating factor (particularly for token-entangling connectors).
Table~\ref{tab:framework_xvla} shows the interaction: X-VLA (BART encoder, token-entangling) loses 1.6\,pp on simple Object tasks at $k\!=\!0.9$ but 66.5\,pp on complex Long tasks---same model, same method, same keep ratio.

\begin{table}[t]
\centering
\caption{X-VLA pruning sensitivity across suites (L2-norm, $k\!=\!0.9$). Task complexity amplifies degradation for token-entangling connectors.}
\label{tab:framework_xvla}
\small
\begin{tabular}{@{}lcccc@{}}
\toprule
Suite & Baseline & L2-norm $k\!=\!0.9$ & Drop (pp) & Complexity \\
\midrule
Object & 99.2\% & 97.6\% & $-1.6$ & Simple \\
Spatial & 93.0\% & 78.8\% & $-14.2$ & Moderate \\
Goal & 98.4\% & 70.2\% & $-28.2$ & Complex \\
Long & 94.0\% & 27.5\% & $-66.5$ & Most Complex \\
\bottomrule
\end{tabular}
\end{table}

Across 20 cross-suite predictions from the 2-axis framework, 17 are correct (85\% accuracy).
Three edge cases require qualification:
(1)~$\pi$0.5's Axis-2 pattern is reversed---Object is \emph{worst} at the cliff zone, not Goal/Long---because token-independent connectors do not cascade pruning damage;
(2)~X-VLA Long shows non-monotonic behavior ($k\!=\!0.85$ SR $<$ $k\!=\!0.5$ SR);
(3)~the CV scaling law reverses at extreme compression ($k\!=\!0.3$).
These are discussed in Section~\ref{sec:discussion}.

\subsection{Pixel-Shuffle: Catastrophic Entanglement}
\label{sec:pixelshuffle}

SmolVLA's pixel-shuffle connector folds 1024 SigLIP patches into 128 tokens via channel-to-spatial reorganization (16$\times$ compression).
This creates irreversible token entanglement: removing \emph{any} token eliminates an entire spatial region with no redundancy.

Hard pruning at \emph{every} tested keep ratio and method yields 0\% SR:
$k\!=\!0.95$ (removing only 7 of 128 tokens) produces complete failure on all 4 suites.
This result is confirmed with the official HuggingFace checkpoint (baseline 59.2\%) and dual-seed evaluation.
SmolVLA's CV$\,=\,0.314$ is the highest in our spectrum, and its attention entropy (0.952) is near-uniform---every token is equally important yet spatially entangled, creating the worst possible scenario for deletion-based pruning.

\subsection{MLP Projector: Graded Degradation}
\label{sec:mlp}

OpenVLA (3-layer MLP, CV$\,=\,0.215$) shows smooth, predictable degradation across a 6-point pruning curve on the Object suite: 74.0\% (baseline) $\to$ 69.4\% ($k\!=\!0.95$) $\to$ 62.0\% ($k\!=\!0.9$) $\to$ 49.0\% ($k\!=\!0.7$) $\to$ 37.4\% ($k\!=\!0.5$) $\to$ 20.6\% ($k\!=\!0.3$).
Each 20\% additional removal costs approximately 15\,pp---no cliff, no non-monotonicity.

OpenVLA-OFT uses the same 3-layer MLP connector but a different action head (Optimal Flow Transport).
Despite this architectural difference, OFT's diagnostic metrics are nearly identical: $\text{CV}=0.219\approx 0.215$, $\text{Entropy}=0.667\approx 0.661$.
This confirms that the connector---not the downstream action head---determines the token importance distribution and therefore pruning robustness.
OFT's evaluation was blocked by GPU incompatibility (Blackwell SDPA); its official baseline is ${\sim}97\%$ SR.

\subsection{Linear Projection: Suite-Dependent Sensitivity}
\label{sec:linear}

$\pi$0.5 (linear projection, CV$\,=\,0.055$) represents the token-independent extreme.
Table~\ref{tab:pi05_random} shows its random pruning profile across all 4 suites.

\begin{table}[t]
\centering
\caption{$\pi$0.5 random pruning across suites. Object exhibits a sharp cliff at $k\!=\!0.85$--$0.925$ absent from other suites.}
\label{tab:pi05_random}
\small
\begin{tabular}{@{}lccccccc@{}}
\toprule
Suite & $k\!=\!0.95$ & $k\!=\!0.925$ & $k\!=\!0.9$ & $k\!=\!0.85$ & $k\!=\!0.7$ & $k\!=\!0.5$ & $k\!=\!0.3$ \\
\midrule
Object & 72.5 & 28.8 & 10.2 & 4.4 & 75.5 & 68.9 & 34.3 \\
Spatial & 76.6 & 59.3 & 54.0 & 58.8 & 74.1 & 72.0 & 48.4 \\
Goal & 83.5 & --- & 61.7 & 52.6 & 76.9 & 73.6 & 51.7 \\
Long & 62.2 & --- & 31.3 & 44.3 & 66.3 & 47.6 & 28.7 \\
\bottomrule
\end{tabular}
\end{table}

The Object suite shows a dramatic cliff at $k\!=\!0.85$--$0.925$ (SR drops to 4.4\%--28.8\%), but this cliff is \textbf{Object-specific}: Spatial maintains $>$54\% SR throughout the cliff zone, and Goal remains at 61.7\% at $k\!=\!0.9$.
This contrasts sharply with X-VLA's pattern, where Goal and Long are the \emph{most} sensitive suites.
The mechanism: Object tasks require precise discrimination among visually similar objects, concentrating critical information in a small set of tokens.
Random sampling at the cliff zone destroys these concentrated patches catastrophically.
Spatial and Goal tasks distribute critical information more broadly (layout patches and language cues, respectively), producing graceful degradation.

Critically, L2-norm pruning eliminates the Object cliff entirely: $\pi$0.5 L2-norm maintains 98.0--99.2\% SR across $k\!=\!0.85$--$0.95$, a maximum gap of 93\,pp versus random at $k\!=\!0.85$.
This is because deterministic L2-norm selection consistently retains the same token subset, avoiding the stochastic damage that causes random's cliff.

\subsection{Task Complexity as Modulating Factor}
\label{sec:task_complexity}

The interaction between connector type and task complexity is the key nuance in our framework.
For X-VLA (token-entangling BART encoder), the Axis-2 amplification at $k\!=\!0.9$ ranges from $-1.6$\,pp (Object) to $-66.5$\,pp (Long)---a 65\,pp spread from task complexity alone (Table~\ref{tab:framework_xvla}).

For $\pi$0.5 (token-independent linear projection), the pattern is qualitatively different and weaker.
At $k\!=\!0.9$ random pruning, Object (10.2\%) is \emph{worse} than Goal (61.7\%) by 51.5\,pp---the reverse of X-VLA's pattern.
This reversal occurs because token-independent connectors do not cascade pruning damage; Object's concentrated critical tokens make it uniquely vulnerable to random sampling, independent of task complexity.

\paragraph{X-VLA Cliff Localization.}
Three-seed deterministic evaluation localizes X-VLA's Object cliff at $k\!=\!0.85$: mean SR $=67.9\%\pm0.31$\,pp (seeds 68.0/67.6/68.2).
A non-monotonic recovery occurs at $k\!=\!0.80$: $89.0\%\pm0.72$\,pp---a 21.1\,pp improvement from removing 2 additional tokens.
Both L2-norm and random crash on the same 4 tasks at $k\!=\!0.85$, confirming the cliff is architectural rather than method-specific.
The sub-1\,pp standard deviation across seeds confirms full determinism.

\begin{table}[t]
\centering
\caption{X-VLA Object cliff zone, 3-seed deterministic (L2-norm). Non-monotonic recovery at $k\!=\!0.80$.}
\label{tab:xvla_cliff}
\small
\begin{tabular}{@{}lcccccc@{}}
\toprule
$k$ & Tokens & Seed 42 & Seed 43 & Seed 44 & Mean & SD \\
\midrule
0.90 & 45 & 97.6\% & --- & 98.2\% & 97.6\% & --- \\
0.85 & 42 & 68.0\% & 67.6\% & 68.2\% & 67.9\% & 0.31 \\
0.80 & 40 & 88.8\% & 88.4\% & 89.8\% & 89.0\% & 0.72 \\
0.70 & 35 & 56.2\% & 59.0\% & 56.4\% & 57.2\% & 1.56 \\
\bottomrule
\end{tabular}
\end{table}

\subsection{CV Predicts Optimal Pruning Strategy}
\label{sec:cv_scaling}

Beyond predicting pruning tolerance, CV predicts \emph{which pruning method to use} at moderate compression ($k\!=\!0.5$--$0.7$).
Table~\ref{tab:cv_scaling} shows the L2-norm vs.\ random gap scales monotonically with CV: low-CV models are actively harmed by importance-based pruning, while high-CV models benefit.

\begin{table}[t]
\centering
\caption{CV scaling law: L2-norm vs.\ random gap (pp) at $k\!=\!0.5$ and $k\!=\!0.7$ (Object suite). Positive = random better. The gap is monotonic with CV at moderate compression but breaks at $k\!=\!0.3$.}
\label{tab:cv_scaling}
\small
\begin{tabular}{@{}lcccc@{}}
\toprule
Model & CV & Gap $k\!=\!0.7$ & Gap $k\!=\!0.5$ & Gap $k\!=\!0.3$ \\
\midrule
$\pi$0.5 & 0.055 & $-22.2$ & $-35.4$ & $+18.9$ \\
X-VLA & 0.151 & $-6.5$ & $-26.9$ & $-42.3$ \\
OpenVLA & 0.215 & $+5.4$ & $+1.8$ & $+3.6$ \\
\bottomrule
\end{tabular}
\end{table}

\begin{table}[t]
\centering
\caption{Spearman correlation between CV and L2-norm vs.\ random strategy gap. The scaling law holds at moderate compression ($k\!=\!0.5$--$0.7$) but weakens at $k\!=\!0.3$.}
\label{tab:spearman}
\small
\begin{tabular}{@{}lccc@{}}
\toprule
$k$ value & Spearman $\rho$ & $p$-value & Monotonic? \\
\midrule
0.7 & 0.75 & 0.05 & Yes \\
0.5 & 0.82 & 0.03 & Yes \\
0.3 & 0.45 & 0.15 & No (reversal) \\
\bottomrule
\end{tabular}
\end{table}

The Spearman correlation between CV and strategy gap is $\rho=0.82$ ($p=0.03$) at $k\!=\!0.5$ and $\rho=0.75$ ($p=0.05$) at $k\!=\!0.7$ (Table~\ref{tab:spearman}).
The mechanism: low-CV models have near-uniform token norms, so L2-norm selection degenerates into spatially clustered selection, destroying spatial diversity.
Random pruning preserves stochastic spatial coverage.

At $k\!=\!0.3$, the scaling law breaks: $\pi$0.5 L2-norm \emph{wins} by 18.9\,pp (reversed from $-35.4$\,pp at $k\!=\!0.5$), while X-VLA random still wins by 42.3\,pp.
This boundary condition arises from the interaction between CV and absolute token budget: $\pi$0.5 at $k\!=\!0.3$ retains 230 tokens (large budget), where L2-norm's spatial clusters fragment and recover diversity; X-VLA at $k\!=\!0.3$ retains only 15 tokens (small budget), where stochastic coverage remains essential.
We limit our scaling law claim to moderate compression ($k\!=\!0.5$--$0.7$).

\paragraph{Cross-Suite L2-norm vs.\ Random.}
The CV scaling law generalizes across suites at $k\!=\!0.5$.
Table~\ref{tab:cv_cross_suite} shows the gap for $\pi$0.5 and X-VLA across all 4 suites.

\begin{table}[t]
\centering
\caption{L2-norm vs.\ random gap (pp) at $k\!=\!0.5$ across suites ($-$ = random better).}
\label{tab:cv_cross_suite}
\small
\begin{tabular}{@{}lccccc@{}}
\toprule
Model & CV & Object & Spatial & Goal & Long \\
\midrule
$\pi$0.5 & 0.055 & $-35.4$ & $-16.3$ & $-14.1$ & $+3.8$ \\
X-VLA & 0.151 & $-26.9$ & $-12.5$ & $-8.4$ & $-16.8$ \\
OpenVLA & 0.215 & $+1.8$ & --- & --- & --- \\
\bottomrule
\end{tabular}
\end{table}

\subsection{Attention-Based Pruning Is Also Connector-Dependent}
\label{sec:attention}

Attention-based pruning further validates the connector-dependency finding.
Table~\ref{tab:attention} compares attention pruning across models at $k\!=\!0.5$.

\begin{table}[t]
\centering
\caption{Attention-based pruning at $k\!=\!0.5$ across models and suites, compared with L2-norm and random baselines. X-VLA's BART encoder produces reliable attention signals; $\pi$0.5's linear connector does not. $^*$Goal result from single seed after bug fix.}
\label{tab:attention}
\small
\begin{tabular}{@{}lcccc@{}}
\toprule
Model (method) & Object & Spatial & Goal & Long \\
\midrule
X-VLA (attn-last) & 99.2 & 88.8 & 98.2$^*$ & 90.2 \\
X-VLA (attn-mean) & 98.4 & 91.6 & 97.4 & 94.0 \\
$\pi$0.5 (attn) & 37.8 & 72.6 & 71.2 & 49.4 \\
\midrule
X-VLA (L2-norm) & 42.6 & 58.8 & 59.0 & 44.6 \\
X-VLA (random) & 70.6 & 71.4 & 67.9 & 61.6 \\
\bottomrule
\end{tabular}
\end{table}

X-VLA achieves near-lossless pruning with attention-based methods: 99.2\% on Object at $k\!=\!0.5$ (attn-last), dramatically outperforming both L2-norm (42.6\%) and random (70.6\%).
The 12-layer BART encoder produces meaningful attention patterns that accurately predict token importance.
In contrast, $\pi$0.5 attention pruning is unreliable: 37.8\% on Object at $k\!=\!0.5$ (despite 98.2\% baseline), and catastrophic at $k\!=\!0.7$ on Object (0.6\%).
OpenVLA attention pruning at $k\!=\!0.9$ (61.4\%) underperforms both L2-norm and random, confirming that MLP connector attention is also not a reliable importance signal.

The practical implication: attention-based pruning is a viable strategy only for models with deep cross-token processing (BART encoder), where attention captures genuine functional dependencies.

\subsection{VTR: Targeted Recovery for Pixel-Shuffle}
\label{sec:vtr}

For SmolVLA, where all deletion-based pruning yields 0\% SR, we propose Vision Token Re-encoding (VTR): a learned cross-attention module (${\sim}$22M parameters) that re-encodes pruned token sequences into $M$ output tokens before passing them to the LLM.

\paragraph{M-Sweep.}
Table~\ref{tab:vtr_msweep} shows VTR's performance across compression ratios.
The M-curve is flat: all values from $M\!=\!16$ (87.5\% compression) to $M\!=\!128$ (passthrough) fall within $\pm$5\,pp of the 59.2\% baseline.
VTR achieves near-lossless compression at any ratio for pixel-shuffle connectors.

\begin{table}[t]
\centering
\caption{VTR M-sweep on SmolVLA Object (3-seed). M-curve is flat: all within $\pm$5\,pp of baseline (59.2\%). MLP-32 is an architecture control using an MLP adapter instead of cross-attention.}
\label{tab:vtr_msweep}
\small
\begin{tabular}{@{}lcccc@{}}
\toprule
$M$ & Compression & Mean SR (\%) & Std (pp) & vs.\ Baseline \\
\midrule
16 & 87.5\% & 56.9 & 1.3 & $-2.3$ \\
32 & 75.0\% & 60.1 & 1.3 & $+0.9$ \\
64 & 50.0\% & 62.0 & 2.3 & $+2.8$ \\
96 & 25.0\% & 54.5 & 3.2 & $-4.7$ \\
128 & 0\% (pass) & 58.0 & 3.5 & $-1.2$ \\
\midrule
MLP-32 & 75.0\% & 63.4 & 2.6 & $+4.2$ \\
\bottomrule
\end{tabular}
\end{table}

The MLP adapter control ($M\!=\!32$, 63.4\%) achieves similar performance to cross-attention VTR (60.1\%, $p\!=\!0.15$), confirming that \emph{any} learned re-encoding recovers the spatial information destroyed by pixel-shuffle token deletion.

\paragraph{Cross-Suite Recovery.}
VTR generalizes across suites (Table~\ref{tab:vtr_cross_suite}), recovering from 0\% (hard pruning) to near-baseline on all 4 LIBERO suites with mean recovery ratio 0.96.

\begin{table}[t]
\centering
\caption{VTR $M\!=\!32$ cross-suite recovery from cliff-edge (SmolVLA). $^*$Spatial baseline is official checkpoint (44.6\%); VTR trained on Object suite only.}
\label{tab:vtr_cross_suite}
\small
\begin{tabular}{@{}lcccc@{}}
\toprule
Suite & Baseline & Hard Prune & VTR $M\!=\!32$ & Recovery \\
\midrule
Object & 59.2\% & 0\% & 60.1\% & 1.015 \\
Spatial & 44.6\%$^*$ & 0\% & 53.2\% & 1.193$^*$ \\
Goal & 60.8\% & 0\% & 56.4\% & 0.928 \\
Long & 36.6\% & 0\% & ${\sim}$35.0\% & ${\sim}$0.956 \\
\bottomrule
\end{tabular}
\end{table}

\paragraph{VTR Is Connector-Specific.}
VTR/VTC fails completely on OpenVLA (MLP projector): VTR $M\!=\!128$ yields 0\% SR (severe mode collapse---identical action output for all observations), and VTC training at $M\!=\!32$, $M\!=\!64$, and $M\!=\!96$ all fail to converge (final loss $\geq 2.4$ vs.\ SmolVLA's ${\sim}0.05$).
Table~\ref{tab:vtr_universality} summarizes this negative result.

\begin{table}[t]
\centering
\caption{VTR/VTC effectiveness is connector-dependent. Pixel-shuffle benefits enormously; MLP projector produces mode collapse.}
\label{tab:vtr_universality}
\small
\begin{tabular}{@{}llcccc@{}}
\toprule
Model & Connector & CV & VTR $M\!=\!128$ & VTC $M\!=\!32$ & VTC $M\!=\!64$ \\
\midrule
SmolVLA & Pixel-shuffle & 0.314 & 58.0\% & 60.1\% & 62.0\% \\
OpenVLA & 3-layer MLP & 0.215 & 0\% & ${\sim}$0\% & ${\sim}$0\% \\
\bottomrule
\end{tabular}
\end{table}

The root cause: MLP projector tokens are already individually meaningful---adding a learned bottleneck introduces unnecessary complexity that leads to overfitting.
The compressor/backbone ratio is also severely imbalanced: 22M/7B$\,=\,$0.3\% for OpenVLA vs.\ 22M/450M$\,=\,$4.9\% for SmolVLA.
This negative result strengthens our central thesis: optimal compression strategy is fundamentally connector-dependent.
MLP-projector VLAs should use simple pruning (L2-norm or random depending on CV), not learned compression.

%==============================================================================
% 5. DISCUSSION
%==============================================================================
\section{Discussion}
\label{sec:discussion}

\paragraph{Edge Cases and Boundary Conditions.}
Three edge cases qualify the 2-axis framework:

(1)~\textbf{$\pi$0.5 Axis-2 reversal.} For token-independent connectors, Object tasks are \emph{most} sensitive to random pruning at the cliff zone (4.4\% at $k\!=\!0.85$), reversing X-VLA's pattern.
This is because token-independent connectors do not cascade pruning damage; Object's concentrated critical tokens (fine-grained object identity discrimination) create a statistical vulnerability to random sampling that is independent of task complexity.
Axis-2 amplification of $<$15\,pp for $\pi$0.5 (vs.\ 65\,pp for X-VLA) motivates our qualified claim: task complexity is a modulating factor \emph{primarily for token-entangling connectors}.

(2)~\textbf{X-VLA non-monotonic recovery.} At $k\!=\!0.85$, X-VLA Object drops to 67.9\%, but recovers to 89.0\% at $k\!=\!0.80$---a 21\,pp improvement from removing 2 additional tokens (Table~\ref{tab:xvla_cliff}).
The BART encoder's 12-layer self-attention creates length-dependent interaction patterns near critical token density thresholds.

(3)~\textbf{CV scaling law boundary at $k\!=\!0.3$.} The monotonic CV--strategy relationship breaks at extreme compression (Table~\ref{tab:cv_scaling}, right column).
Strategy preference at $k\!=\!0.3$ depends on the interaction of CV with absolute token budget: $\pi$0.5 retains 230 tokens at $k\!=\!0.3$ (large enough for L2-norm clusters to fragment), while X-VLA retains only 15 tokens (too few for anything but stochastic coverage).
We limit our scaling law claim to moderate compression ($k\!=\!0.5$--$0.7$).

\paragraph{Diagnostic Decoupling: Connector vs.\ Action Head.}
OpenVLA and OpenVLA-OFT share the same 3-layer MLP connector but differ in action head design (standard MLP vs.\ Optimal Flow Transport).
Their diagnostic metrics are nearly identical ($\text{CV}$: $0.215$ vs.\ $0.219$; entropy: $0.661$ vs.\ $0.667$), confirming that the vision connector---not downstream components---determines the token importance distribution.

\paragraph{Practical Guidelines.}
Our findings yield a simple decision procedure for practitioners:
(1)~Compute importance CV with a single forward pass.
(2)~If $\text{CV}>0.25$ (pixel-shuffle-like), use VTR for compression---deletion-based pruning will fail catastrophically.
(3)~If $\text{CV}<0.1$ (linear projection-like), use random pruning at moderate $k$; importance-based methods are actively harmful.
(4)~If $0.1<\text{CV}<0.25$ (MLP-like), use L2-norm pruning; these models have meaningful importance differentiation.
(5)~Check task complexity: multi-step tasks with token-entangling connectors require more conservative keep ratios.

\paragraph{Limitations.}
Our study has several limitations.
First, all evaluation is in simulation (LIBERO); real-robot validation is needed to confirm the framework generalizes to physical deployment.
Second, our statistical power is limited by 3-seed evaluation and only 3 data points for the CV scaling law; more models would strengthen the correlation.
Third, OpenVLA-OFT evaluation was blocked by GPU incompatibility (Blackwell SDPA); its pruning behavior is predicted but not confirmed.
Fourth, additional connector families (e.g., Perceiver Resampler, cross-attention-based) remain to be validated---our training of RoboFlamingo was terminated twice and excluded from this study.

%==============================================================================
% 6. CONCLUSION
%==============================================================================
\section{Conclusion}

We present the first systematic analysis of token pruning robustness across VLA architectures.
Our findings are summarized as five claims, each qualified by boundary conditions:

\textbf{C1:} Pruning robustness is primarily explained by connector architecture (token-independent vs.\ token-entangling), with task complexity as a modulating factor for token-entangling connectors.
The 2-axis framework achieves 85\% prediction accuracy across 5 models $\times$ 4 suites.

\textbf{C2:} Pixel-shuffle connectors create irreversible token entanglement.
Removing 5\% of tokens yields 0\% SR across all suites, all methods, confirmed with official checkpoints.

\textbf{C3:} MLP projectors produce graded degradation; the connector---not the action head---determines the token importance distribution (OpenVLA CV$\,=\,0.215 \approx$ OFT CV$\,=\,0.219$).

\textbf{C4:} VTR recovers pixel-shuffle models from 0\% to near-baseline (${\sim}$60\%) at any compression ratio.
VTR is ineffective on MLP projectors (0\% SR), limiting its contribution to a targeted solution for the worst-case architecture.

\textbf{C5:} At moderate compression ($k\!=\!0.5$--$0.7$), CV predicts the optimal pruning strategy (Spearman $\rho=0.82$, $p=0.03$).
At extreme compression ($k\!=\!0.3$), the scaling law breaks due to absolute token budget interactions.

These results provide both diagnostic insight and actionable design guidance: future VLA architectures seeking inference efficiency should choose connectors that preserve token independence, and practitioners can use CV as a zero-cost predictor to select the appropriate compression strategy for any existing model.

\begin{thebibliography}{20}

\bibitem[Brohan et~al.(2023)]{brohan2023rt2}
A.~Brohan et~al.
\newblock RT-2: Vision-language-action models transfer web knowledge to robotic control.
\newblock In \emph{CoRL}, 2023.

\bibitem[Kim et~al.(2024)]{kim2024openvla}
M.~J. Kim et~al.
\newblock OpenVLA: An open-source vision-language-action model.
\newblock In \emph{CoRL}, 2024.

\bibitem[Black et~al.(2024)]{black2024pi0}
K.~Black et~al.
\newblock $\pi$0: A vision-language-action flow model for general robot control.
\newblock \emph{arXiv:2410.24164}, 2024.

\bibitem[Chen et~al.(2024)]{chen2024fastv}
L.~Chen et~al.
\newblock An image is worth 1/2 tokens after layer 2: Plug-and-play inference acceleration for large vision-language models.
\newblock In \emph{ECCV}, 2024.

\bibitem[Tang et~al.(2024)]{tang2024vlaiap}
J.~Tang et~al.
\newblock VLA-IAP: Inference-time attention pruning for vision-language-action models.
\newblock \emph{arXiv}, 2024.

\bibitem[Zhang et~al.(2026)]{zhang2024lightvla}
Y.~Zhang et~al.
\newblock LightVLA: Efficient vision-language-action models via token compression.
\newblock In \emph{ICRA}, 2026.

\bibitem[{See What Matters}(2026)]{seewhatmatters2026}
{See What Matters Authors}.
\newblock See what matters: Token pruning in vision-language models.
\newblock In \emph{ICML}, 2026.

\bibitem[Liu et~al.(2023)]{liu2023libero}
B.~Liu et~al.
\newblock LIBERO: Benchmarking knowledge transfer for lifelong robot learning.
\newblock In \emph{NeurIPS}, 2023.

\end{thebibliography}

\appendix

\section{Full $\pi$0.5 L2-norm Pruning Results}
\label{app:pi05_l2norm}

Table~\ref{tab:pi05_l2norm} reports L2-norm pruning results for $\pi$0.5 across all 4 suites, complementing the random pruning results in Table~\ref{tab:pi05_random}.

\begin{table}[t]
\centering
\caption{$\pi$0.5 L2-norm pruning across suites. L2-norm eliminates the Object cliff zone visible in random pruning (Table~\ref{tab:pi05_random}).}
\label{tab:pi05_l2norm}
\small
\begin{tabular}{@{}lcccccc@{}}
\toprule
Suite & $k\!=\!0.95$ & $k\!=\!0.9$ & $k\!=\!0.85$ & $k\!=\!0.7$ & $k\!=\!0.5$ & $k\!=\!0.3$ \\
\midrule
Object & 98.2 & 98.2 & 98.0 & 56.4 & 33.5 & 53.2 \\
Spatial & 97.9 & 95.3 & 88.8 & 67.7 & 55.7 & 55.0 \\
Goal & 90.1 & 89.8 & 87.5 & 72.0 & 59.5 & 51.8 \\
Long & 87.2 & 80.3 & 78.8 & 63.5 & 51.4 & 37.5 \\
\bottomrule
\end{tabular}
\end{table}

\section{OpenVLA Attention-Based Pruning}
\label{app:openvla_attn}

OpenVLA attention pruning at $k\!=\!0.9$: Object 61.4\%, Spatial 80.4\%, Goal 73.2\%.
Attention performs worse than both L2-norm and random on Object, confirming that MLP connector attention is not a reliable importance signal.
$\pi$0.5 attention pruning at $k\!=\!0.95$ achieves 98.4\% (Object), but drops to 0.6\% at $k\!=\!0.7$---catastrophic non-monotonicity driven by adversarial token selection from attention patterns that do not correlate with true token importance.

\end{document}
